% Combined group presentation.
% Order follows the outline:
%   0. Title page
%   1. Outline, 2. Goal                  (Outline_Goal.tex)
%   3. Architecture                       (presentation_kurz.tex, restyled)
%   4. Literature                         (TO BE ADDED, see marker below)
%   5. Parameters, 6. Experiments        (Philip.tex, restyled)
%   7. Discussion, 8. Outlook, closing   (Hadi_Discussion_Outlook.tex)
\documentclass[aspectratio=169]{beamer}
\usepackage[scaled=0.98]{helvet}
\renewcommand{\familydefault}{\sfdefault}
\usepackage{booktabs}
\usepackage{array}
\usepackage{multirow}
\usepackage{graphicx}
\usepackage{amsmath,amssymb}
\usepackage{tikz}
\usetikzlibrary{arrows.meta,positioning}
\usepackage{xcolor}

% Reserved academic palette with a title strip and no shaded callout boxes.
\definecolor{ink}{HTML}{26343D}
\definecolor{muted}{HTML}{59666C}
\definecolor{accent}{HTML}{8B6A45}
\definecolor{header}{HTML}{273C3D}
\definecolor{paper}{HTML}{FBFBF9}
\setbeamercolor{background canvas}{bg=paper}
\setbeamercolor{normal text}{fg=ink}
\setbeamercolor{frametitle}{fg=white,bg=header}
\setbeamercolor{structure}{fg=accent}
\setbeamercolor{footline}{fg=muted}
\setbeamerfont{frametitle}{size=\Large,series=\bfseries}
\setbeamerfont{footline}{size=\scriptsize}
\setbeamertemplate{navigation symbols}{}
\setbeamertemplate{frametitle}{\nointerlineskip\begin{beamercolorbox}[wd=\paperwidth,ht=0.72cm,dp=0.15cm,leftskip=0.62cm]{frametitle}\usebeamerfont{frametitle}\insertframetitle\end{beamercolorbox}}
\setbeamertemplate{footline}{\hfill\usebeamerfont{footline}\insertframenumber/\inserttotalframenumber\hspace{0.55cm}\vspace{0.25cm}}
\setbeamertemplate{itemize item}{\textcolor{accent}{\raisebox{0.15ex}{\scriptsize$\bullet$}}}
\setbeamertemplate{itemize subitem}{\textcolor{accent}{$\bullet$}}
\setbeamertemplate{enumerate item}{\textcolor{accent}{\textbf{\insertenumlabel}}}
\setlength{\tabcolsep}{8pt}

% Helpers.
\newcolumntype{L}[1]{>{\raggedright\arraybackslash}p{#1}}
\newcommand{\outlineitem}[2]{%
  \textcolor{accent}{\textbf{#1}}\hspace{0.8em}\textbf{#2}\par\vspace{0.55em}}
\newcommand{\outlinesub}[1]{%
  \hspace{2.1em}\textcolor{muted}{\small #1}\par\vspace{0.25em}}

\begin{document}

%######################################################################
% TITLE PAGE (shown during the introduction)
%######################################################################

\begingroup
\usebackgroundtemplate{\vbox to \paperheight{\hsize=\paperwidth\nointerlineskip
  \hbox to \paperwidth{\color{header}\vrule width \paperwidth height 5.0cm depth 0pt}\nointerlineskip
  \hbox to \paperwidth{\color{accent}\vrule width \paperwidth height 0.1cm depth 0pt}\vskip 0.4cm
  \hbox to \paperwidth{\hfil\IfFileExists{title_image.png}{\includegraphics[width=5.0cm]{title_image.png}}{}\hspace{0.8cm}}\vfil}}
\begin{frame}[plain,t]
\vspace{0.85cm}
\hspace{0.45cm}\begin{minipage}{0.84\paperwidth}
{\color{white}\fontsize{27}{32}\selectfont\bfseries Time Series Forecasting\\in the Frequency Domain\par}
\vspace{0.45cm}
{\color{white}\normalsize Where periodic modeling helps lightweight forecasters, and where it fails\par}
\end{minipage}\par
\vspace{1.55cm}
\hspace{0.45cm}\begin{minipage}{0.84\paperwidth}
\small
\textbf{Aaditya Chaturvedi \quad Abdelhadi Hadj Sahraoui}\\[0.15em]
\textbf{Philip Gentner \quad Mohamad Taleb}\par\vspace{0.3em}
\textcolor{muted}{\footnotesize Supervisor: Jung Min Choi}\par\vspace{0.45em}
\textbf{SRP interim presentation}\par\vspace{0.3em}
\textcolor{muted}{\footnotesize Information Systems and Machine Learning Lab (ISMLL)\\
University of Hildesheim}
\end{minipage}
\end{frame}
\endgroup

%######################################################################
% PART 1 --- OUTLINE AND GOAL
%######################################################################

\begin{frame}{Outline}
\normalsize
\vspace{0.6em}
\outlineitem{1}{What is the goal?}
\outlineitem{2}{Our architecture}
\outlineitem{3}{Literature}
\outlinesub{CycleNet \quad$\cdot$\quad HDMixer \quad$\cdot$\quad PatchTST}
\vspace{0.3em}
\outlineitem{4}{Which parameters do what?}
\outlineitem{5}{Experiments}
\outlineitem{6}{Discussion: why did something work or not work?}
\outlineitem{7}{Outlook: what do we want to do next?}
\end{frame}

\begin{frame}{What is the goal?}
\small
\textit{Find out when periodic modeling helps lightweight forecasters, and when it does not.}\par
\vspace{0.55em}
{\setlength{\fboxsep}{7pt}\setlength{\fboxrule}{0.45pt}%
\noindent\color{black}\fbox{\begin{minipage}{\dimexpr\linewidth-2\fboxsep-2\fboxrule\relax}
\textbf{Research question}\par\vspace{0.25em}
When does frequency-domain and periodic modeling actually help versus fail in lightweight multivariate time series forecasting?
\end{minipage}}}
\vspace{0.75em}
\begin{itemize}
  \setlength{\itemsep}{0.65em}
  \item \textbf{Starting point:} CycleNet stores a learnable cyclic pattern, subtracts it from the input, forecasts the remaining residual with a small backbone, and adds the pattern back.
  \item \textbf{Open limitation:} The cycle length is chosen by hand for each dataset (for example 24, 96, 144 or 168). The paper itself lists varying cycle lengths across channels as unresolved.
  \item \textbf{Our approach:} Replace these manual choices with decisions made from the data, and log every decision so that a result can be traced back to a property of the dataset.
\end{itemize}
\end{frame}

\begin{frame}{What is the goal? Data-driven settings}
\small
\textit{The forward pass stays the same. What changes is how its settings are chosen.}\par
\vspace{0.7em}
\renewcommand{\arraystretch}{1.08}
\begin{tabular}{@{}>{\raggedright\arraybackslash}p{3.0cm}>{\raggedright\arraybackslash}p{3.4cm}>{\raggedright\arraybackslash}p{6.2cm}@{}}
\toprule
\textbf{Component} & \textbf{Original CycleNet} & \textbf{Our version} \\
\midrule
Cycle length & Set by hand per dataset & FFT peaks on the training series, kept only if the autocorrelation exceeds its 95\% noise bound \\
Number of cycles & One & One to three; an extra period is added only if it explains at least 2\% more variance \\
Channel handling & Channel-independent & A correlation check decides; a mixing layer is added only for strongly correlated channels \\
Backbone & Linear or MLP & Linear, MLP, CNN or spectral-conditioned patch mixer (SOR) \\
\bottomrule
\end{tabular}
\vspace{0.5em}

{\footnotesize\textcolor{muted}{Current limitation: periods are detected on the channel-averaged training series, so all channels share the same cycles for now.}}
\end{frame}

\begin{frame}{What is the goal? What we want to find out}
\normalsize
{\small\textit{Each question is answered with the same fixed protocol on every dataset.}}\par
\vspace{0.8em}
\textcolor{accent}{\textbf{1}}\quad \textbf{Can detected cycles replace hand-set ones?}
\par\vspace{0.3em}
{\small Compare the detected periods with the standard values and compare the resulting forecasts.}

\vspace{0.5em}
\textcolor{accent}{\textbf{2}}\quad \textbf{When do extra components help?}
\par\vspace{0.3em}
{\small Check whether multiple cycles and channel mixing pay off only on the datasets where the diagnostics suggest it, and report where the gain is small.}

\vspace{0.5em}
\textcolor{accent}{\textbf{3}}\quad \textbf{Which backbone suits the residual, and why?}
\par\vspace{0.3em}
{\small Compare four backbones on the same de-cycled residual, then use ablations to identify what drives any difference.}

\vspace{0.4em}
{\footnotesize\textcolor{muted}{Protocol: lookback 336, chronological splits (12/4/4 months on ETT, 70/10/20 otherwise), scaler fit on the training split only.}}
\end{frame}

%######################################################################
% PART 2 --- OUR ARCHITECTURE
%######################################################################

\begin{frame}{Normalized frequency metrics}
\small
\textit{Raw MSE is scale-dependent, so we use two scale-free diagnostic metrics.}\par
\vspace{0.5em}
\begin{columns}[T]
\column{0.5\textwidth}
\textbf{1. Fraction of variance unexplained, $FVU_{\text{band}}$}
\[
  FVU_{\text{band}} = \frac{\sum_{k \in \mathcal{B}} w_k |E(k)|^2}{\sum_{k \in \mathcal{B}} w_k |Y(k) - \bar{Y}(k)|^2}
\]
\begin{itemize}
  \setlength{\itemsep}{0.35em}
  \item $FVU < 1.0$: the model learns genuine dynamics in that band.
  \item $FVU \approx 1.0$: no better than predicting a flat mean line.
  \item $FVU > 1.0$: the model hallucinates power and amplifies noise.
\end{itemize}
\column{0.5\textwidth}
\textbf{2. Log-spectral distance, LSD (in dB)}
\[
  \text{LSD} = \sqrt{\frac{1}{M} \sum_{k=0}^{M-1} \left( 10 \log_{10} \frac{|\hat{Y}(k)|^2 + \epsilon}{|Y(k)|^2 + \epsilon} \right)^2}
\]
\begin{itemize}
  \setlength{\itemsep}{0.35em}
  \item Measures distortion of the power spectral density over the full Nyquist range.
  \item Directly penalizes harmonic damping and spectral collapse.
\end{itemize}
\end{columns}
\end{frame}

\begin{frame}{Our architecture: CycleNet-Pro}
\small
\textit{Explicit periodic memory plus modular residual backbones.}\par
\vspace{0.2em}
\begin{columns}[T]
\column{0.54\textwidth}
\begin{enumerate}
  \setlength{\itemsep}{0.4em}
  \item \textbf{RevIN instance normalization:} guards against global non-stationary distribution shifts.
  \item \textbf{Automated period discovery:} extracts multiple cycles via FFT and ACF, with no hardcoded periods.
  \item \textbf{Recurrent cyclic state queues:} explicit learnable periodic parameters $\mathbf{Q}^{(k)}$.
  \item \textbf{Modular residual backbones:} de-cycled inputs are routed through a specialized backbone.
  \item \textbf{Residual channel mixing:} mixes channels only when they are correlated.
\end{enumerate}
\column{0.46\textwidth}
\centering
\tikzset{flow/.style={draw=ink, rounded corners=3pt, line width=0.5pt, fill=white, text width=4.7cm, minimum height=0.75cm, align=center, font=\footnotesize, inner sep=3pt}}
\begin{tikzpicture}[node distance=0.22cm, >={Stealth[length=2.2mm]}]
  \node[flow, fill=muted!10] (inp) {Input series $\mathbf{X} \in \mathbb{R}^{L \times C}$};
  \node[flow, fill=accent!12, below=of inp] (det) {Automated period engine\\(FFT peaks + ACF filter)};
  \node[flow, fill=accent!12, below=of det] (dec) {De-cycling\\$\mathbf{X}_{\text{res}} = \mathbf{X} - \sum_k \mathbf{Q}^{(k)}$};
  \node[flow, below=of dec] (bb) {Modular residual backbone\\(Linear / MLP / CNN / SCPM)};
  \node[flow, fill=muted!10, below=of bb] (rec) {Recombination and RevIN out\\$\hat{\mathbf{Y}} = \mathbf{Y}_{\text{res}} + \mathbf{Y}_{\text{cycle}}$};
  \foreach \a/\b in {inp/det,det/dec,dec/bb,bb/rec} \draw[->, line width=0.6pt] (\a) -- (\b);
\end{tikzpicture}
\end{columns}
\end{frame}

\begin{frame}{Stage 1: automated period discovery (FFT + ACF)}
\small
\textit{Prior models force a hardcoded period ($p=24$ or $p=96$). CycleNet-Pro finds the periods in the training data.}\par
\vspace{0.7em}
\begin{columns}[T]
\column{0.32\textwidth}
\textcolor{accent}{\textbf{1}}\quad \textbf{FFT peak ranking}\par\vspace{0.3em}
\begin{itemize}
  \setlength{\itemsep}{0.4em}
  \item Power spectrum $P(f) = |\text{rFFT}(\mathbf{x})|^2$.
  \item Candidate peaks from local extrema.
  \item Merge spectral-leakage clusters within a 12\% relative distance.
\end{itemize}
\column{0.34\textwidth}
\textcolor{accent}{\textbf{2}}\quad \textbf{ACF noise filter}\par\vspace{0.3em}
\begin{itemize}
  \setlength{\itemsep}{0.4em}
  \item Validate each candidate with the empirical autocorrelation.
  \item 95\% white-noise threshold: $\text{Bound}_{\text{ACF}} = \dfrac{1.96}{\sqrt{N_{\text{train}}}}$
  \item Rejects spurious spectral peaks.
\end{itemize}
\column{0.34\textwidth}
\textcolor{accent}{\textbf{3}}\quad \textbf{Variance decision}\par\vspace{0.3em}
\begin{itemize}
  \setlength{\itemsep}{0.4em}
  \item Marginal variance reduction of each added cycle: $\Delta\text{Var}(k) \ge 2.0\%$.
  \item Accepts secondary and tertiary cycles (for example intra-day plus weekly).
\end{itemize}
\end{columns}
\end{frame}

\begin{frame}{Stage 2: recurrent periodic queues and de-cycling}
\small
\begin{columns}[T]
\column{0.56\textwidth}
\textbf{Explicit recurrent state}\par\vspace{0.2em}
\begin{itemize}
  \setlength{\itemsep}{0.4em}
  \item For each period $p_k \in \mathcal{P}$, keep a learnable state matrix $\mathbf{Q}^{(k)} \in \mathbb{R}^{p_k \times C}$.
  \item With global time index $\tau$, retrieve the stored pattern by circular phase indexing:
  \[ \mathbf{C}^{(k)}_t = \mathbf{Q}^{(k)}_{(\tau + t) \bmod p_k}, \quad t \in [0, L-1] \]
  \item \textbf{Residual de-cycling:} subtract the periodic components from the RevIN-normalized input:
  \[ \mathbf{X}_{\text{res}} = \mathbf{X}_{\text{norm}} - \sum_{k=1}^{K} \mathbf{C}^{(k)} \]
\end{itemize}
\column{0.44\textwidth}
\textbf{Why de-cycling works}\par\vspace{0.2em}
\begin{itemize}
  \setlength{\itemsep}{0.4em}
  \item The residual holds only non-periodic dynamics: trends, transient shocks, regime changes.
  \item The backbone does not waste capacity on repetitive daily cycles.
  \item Periodic extrapolation to horizon $H$ is exact:
  \[ \mathbf{Y}_{\text{cycle},t} = \sum_{k=1}^{K} \mathbf{Q}^{(k)}_{(\tau + L + t) \bmod p_k} \]
\end{itemize}
\end{columns}
\end{frame}

\begin{frame}{Modular residual backbones}
\footnotesize
\textit{Four interchangeable backbones forecast the de-cycled residual.}\par
\vspace{0.3em}
\begin{columns}[T]
\column{0.48\textwidth}
\textcolor{accent}{\textbf{1}}\quad \textbf{Linear and MLP}\par\vspace{0.3em}
\begin{itemize}
  \setlength{\itemsep}{0.3em}
  \item \textbf{Linear:} channel-independent $\mathbf{W} \in \mathbb{R}^{H \times L}$, a parameter-efficient baseline.
  \item \textbf{MLP:} two layers with hidden size 512 and GELU, for non-linear temporal dynamics.
\end{itemize}
\vspace{0.5em}
\textcolor{accent}{\textbf{2}}\quad \textbf{1D-CNN (edge scale)}\par\vspace{0.3em}
\begin{itemize}
  \setlength{\itemsep}{0.3em}
  \item One shared 1D convolution across time and channels, with linear length interpolation ($L \to H$).
  \item Cuts parameters by \textbf{97.7\%} (5.5k versus 241k on ETT).
\end{itemize}
\column{0.52\textwidth}
\textcolor{accent}{\textbf{3}}\quad \textbf{SOR: spectral-conditioned patch mixer}\par\vspace{0.3em}
\begin{itemize}
  \setlength{\itemsep}{0.3em}
  \item \textbf{Intra-patch mixer:} token MLP along time inside each patch, plus a feature MLP.
  \item \textbf{Inter-patch mixer:} neighbor-aware gated MLP across the patch order, plus a feature MLP.
  \item \textbf{FiLM spectral conditioning:} reads orthonormal DCT-II log-energy bands and modulates the patch tokens:
  \[ \mathbf{Z}_{\text{cond}} = \mathbf{Z} \odot (1 + \boldsymbol{\gamma}) + \boldsymbol{\beta} \]
  Zero initialization starts training from the identity.
\end{itemize}
\end{columns}
\end{frame}

\begin{frame}[t]{SOR Architecture: Spectral-Conditioned Patch Mixer}
\centering
\vspace{-0.25cm}
\resizebox{!}{0.82\textheight}{%
    \IfFileExists{sor_presentation.tikz}{% SORMixer Presentation Architecture Diagram (Wide 16:9 Layout)
\providecolor{sorblue}{RGB}{24, 76, 120}
\providecolor{sororange}{RGB}{214, 90, 49}
\providecolor{sorgreen}{RGB}{46, 117, 89}
\providecolor{sordark}{RGB}{45, 52, 54}
\providecolor{sorlight}{RGB}{245, 246, 250}

\begin{tikzpicture}[
    >=stealth,
    font=\small,
    box/.style={draw, rounded corners=3pt, thick, align=center, inner sep=4.5pt, fill=white, font=\small},
    input/.style={box, fill=blue!10, draw=sorblue!80, text width=4.6cm},
    fold/.style={box, fill=cyan!12, draw=cyan!70!black, text width=4.6cm},
    patch/.style={box, fill=teal!12, draw=teal!70!black, text width=4.6cm},
    intra/.style={box, fill=orange!15, draw=sororange!80, text width=4.6cm},
    pool/.style={box, fill=teal!18, draw=teal!80!black, text width=4.6cm},
    film_node/.style={box, fill=purple!20, draw=purple!80!black, text width=4.6cm},
    inter/.style={box, fill=yellow!30!orange!18, draw=sororange!90, text width=4.6cm},
    head/.style={box, fill=cyan!18, draw=cyan!80!black, text width=4.6cm},
    output/.style={box, fill=sorgreen!20, draw=sorgreen!80!black, text width=4.6cm},
    spec/.style={box, fill=violet!12, draw=violet!70!black, text width=4.6cm},
    card/.style={draw=sordark!40, fill=sorlight, rounded corners=3pt, inner sep=5pt, font=\footnotesize, align=left}
]

% =========================================================================
% STAGE 1: Patch Formation & Intra-Patch Modeling (x = 0)
% =========================================================================
\node[font=\bfseries\normalsize, text=sorblue, anchor=center] (s1_title) at (0, 0.4) {1. Patch \& Intra-Mixer};

\node[input, below=0.22cm of s1_title] (in_res) {
    \textbf{Residual Input} $\mathbf{X}_{\text{res}} \in \mathbb{R}^{B \times L \times C}$\\
    \footnotesize De-cycled residual from CycleNet-Pro
};

\node[fold, below=0.22cm of in_res] (fold) {
    \textbf{Channel Fold} $\mathbf{X}_c$\\
    \footnotesize Reshape to $[B \cdot C, L]$ (Channel-Indep.)
};

\node[patch, below=0.22cm of fold] (patching) {
    \textbf{Unfold \& Point Embed}\\
    \footnotesize Patches: $[B \cdot C, N, P] \implies \mathbf{H}_0 \in [B \cdot C, N, P, D]$
};

\node[intra, below=0.22cm of patching] (intra) {
    \textbf{Intra-Patch Mixer}\\
    \footnotesize Time-MLP ($P \to 2P \to P$) + Feat-MLP ($D \to 2D \to D$)
};

\node[pool, below=0.22cm of intra] (pool) {
    \textbf{Patch Token Pooling}\\
    \footnotesize $\text{Linear}(P \cdot D \to D) \implies \mathbf{H} \in [B \cdot C, N, D]$
};

\draw[->, thick, color=sorblue] (in_res) -- (fold);
\draw[->, thick, color=teal!80!black] (fold) -- (patching);
\draw[->, thick, color=teal!80!black] (patching) -- (intra);
\draw[->, thick, color=sororange!90!black] (intra) -- (pool);

% =========================================================================
% STAGE 2: DCT-II Spectral Conditioning Branch (x = 5.5)
% =========================================================================
\node[font=\bfseries\normalsize, text=violet!85!black, anchor=center] (s2_title) at (5.5, 0.4) {2. Spectral FiLM Branch};

\node[spec, below=0.22cm of s2_title] (dct) {
    \textbf{Orthonormal DCT-II}\\
    \footnotesize $\mathbf{C}_{\text{dct}} = \mathbf{X}_c \mathbf{D}^T \in \mathbb{R}^{B \cdot C \times L}$ (Energy compaction)
};

\node[spec, below=0.22cm of dct] (spec_power) {
    \textbf{Log-Power Subbands}\\
    \footnotesize $\mathbf{S} = \log(1 + \mathbf{C}_{\text{dct}}^2) \to \text{Pool}(K=4) \implies \bar{\mathbf{s}} \in [B \cdot C, K]$
};

\node[spec, below=0.22cm of spec_power] (film_mlp) {
    \textbf{Zero-Init FiLM Generator}\\
    \footnotesize MLP: $K \to \text{hidden} \to 2D \implies (\boldsymbol{\gamma}, \boldsymbol{\beta}) = (\mathbf{0}, \mathbf{0})$
};

\node[film_node, below=0.28cm of film_mlp] (film) {
    \textbf{FiLM Spectral Modulation}\\
    $\mathbf{H}_{\text{cond}} = \mathbf{H} \odot (1 + \boldsymbol{\gamma}) + \boldsymbol{\beta}$\\
    \footnotesize \textit{Step 0: Identity Initialization ($\mathbf{H}_{\text{cond}} = \mathbf{H}$)}
};

% Connectors in Stage 2
\draw[->, thick, color=violet!80!black] (dct) -- (spec_power);
\draw[->, thick, color=violet!80!black] (spec_power) -- (film_mlp);
\draw[->, thick, color=purple!90!black] (film_mlp) -- (film)
    node[midway, right, font=\footnotesize] {$(\boldsymbol{\gamma}, \boldsymbol{\beta})$};

% =========================================================================
% STAGE 3: Inter-Patch Mixer & Forecasting Head (x = 11.0)
% =========================================================================
\node[font=\bfseries\normalsize, text=sororange!95!black, anchor=center] (s3_title) at (11.0, 0.4) {3. Inter-Patch \& Forecast};

\node[inter, below=0.22cm of s3_title] (inter) {
    \textbf{Neighbor-Gated Inter-Mixer}\\
    \footnotesize Context: $[\mathbf{h}_{i-1}, \mathbf{h}_i, \mathbf{h}_{i+1}] \in 3D \to \text{MLP}_{3D \to D}$\\
    \footnotesize Dynamic Gate $g_i \in [0, 1]$ + Feat-MLP
};

\node[head, below=0.22cm of inter] (head) {
    \textbf{Forecasting Head}\\
    \footnotesize $\text{Linear}(N \cdot D \to D) \to \text{GELU} \to \text{Linear}(D \to H)$
};

\node[output, below=0.22cm of head] (out_res) {
    \textbf{Residual Forecast} $\hat{\mathbf{Y}}_{\text{res}} \in \mathbb{R}^{B \times H \times C}$\\
    \footnotesize Unfold channels + Periodic forecast $\hat{\mathbf{Y}}_{\text{cycle}}$
};

\node[card, below=0.22cm of out_res, text width=4.6cm] (summary_card) {
    \textbf{Key SORMixer Specifications:}\\
    $\bullet$ \textbf{Params}: $\sim \mathbf{100\text{k}}$ ($10\times$ smaller than PatchTST)\\
    $\bullet$ \textbf{Complexity}: $\mathcal{O}(N)$ strictly linear in patches\\
    $\bullet$ \textbf{Channel}: Independent residual modeling\\
    $\bullet$ \textbf{Conditioning}: Orthonormal DCT-II FiLM
};

% Connectors between Stages
% Branch from Stage 1 Fold to Stage 2 DCT
\draw[->, thick, color=violet!80!black] (fold.east) -- ++(0.45, 0) node[midway, above, font=\footnotesize] {$\mathbf{X}_c$} |- (dct.west);

% Stream from Stage 1 Pool to Stage 2 FiLM
\draw[->, thick, color=teal!80!black] (pool.east) -- ++(0.45, 0) node[midway, above, font=\footnotesize] {$\mathbf{H}$} |- (film.west);

% Stream from Stage 2 FiLM to Stage 3 Inter-Patch (centered in corridor)
\draw[->, thick, color=purple!80!black] (film.east) -- ++(0.45, 0) node[midway, above, font=\footnotesize] {$\mathbf{H}_{\text{cond}}$} |- (inter.west);

\draw[->, thick, color=sororange!90!black] (inter) -- (head);
\draw[->, thick, color=sorgreen!80!black] (head) -- (out_res);

\end{tikzpicture}
}{% SORMixer Presentation Architecture Diagram (Wide 16:9 Layout)
\providecolor{sorblue}{RGB}{24, 76, 120}
\providecolor{sororange}{RGB}{214, 90, 49}
\providecolor{sorgreen}{RGB}{46, 117, 89}
\providecolor{sordark}{RGB}{45, 52, 54}
\providecolor{sorlight}{RGB}{245, 246, 250}

\begin{tikzpicture}[
    >=stealth,
    font=\small,
    box/.style={draw, rounded corners=3pt, thick, align=center, inner sep=4.5pt, fill=white, font=\small},
    input/.style={box, fill=blue!10, draw=sorblue!80, text width=4.6cm},
    fold/.style={box, fill=cyan!12, draw=cyan!70!black, text width=4.6cm},
    patch/.style={box, fill=teal!12, draw=teal!70!black, text width=4.6cm},
    intra/.style={box, fill=orange!15, draw=sororange!80, text width=4.6cm},
    pool/.style={box, fill=teal!18, draw=teal!80!black, text width=4.6cm},
    film_node/.style={box, fill=purple!20, draw=purple!80!black, text width=4.6cm},
    inter/.style={box, fill=yellow!30!orange!18, draw=sororange!90, text width=4.6cm},
    head/.style={box, fill=cyan!18, draw=cyan!80!black, text width=4.6cm},
    output/.style={box, fill=sorgreen!20, draw=sorgreen!80!black, text width=4.6cm},
    spec/.style={box, fill=violet!12, draw=violet!70!black, text width=4.6cm},
    card/.style={draw=sordark!40, fill=sorlight, rounded corners=3pt, inner sep=5pt, font=\footnotesize, align=left}
]

% =========================================================================
% STAGE 1: Patch Formation & Intra-Patch Modeling (x = 0)
% =========================================================================
\node[font=\bfseries\normalsize, text=sorblue, anchor=center] (s1_title) at (0, 0.4) {1. Patch \& Intra-Mixer};

\node[input, below=0.22cm of s1_title] (in_res) {
    \textbf{Residual Input} $\mathbf{X}_{\text{res}} \in \mathbb{R}^{B \times L \times C}$\\
    \footnotesize De-cycled residual from CycleNet-Pro
};

\node[fold, below=0.22cm of in_res] (fold) {
    \textbf{Channel Fold} $\mathbf{X}_c$\\
    \footnotesize Reshape to $[B \cdot C, L]$ (Channel-Indep.)
};

\node[patch, below=0.22cm of fold] (patching) {
    \textbf{Unfold \& Point Embed}\\
    \footnotesize Patches: $[B \cdot C, N, P] \implies \mathbf{H}_0 \in [B \cdot C, N, P, D]$
};

\node[intra, below=0.22cm of patching] (intra) {
    \textbf{Intra-Patch Mixer}\\
    \footnotesize Time-MLP ($P \to 2P \to P$) + Feat-MLP ($D \to 2D \to D$)
};

\node[pool, below=0.22cm of intra] (pool) {
    \textbf{Patch Token Pooling}\\
    \footnotesize $\text{Linear}(P \cdot D \to D) \implies \mathbf{H} \in [B \cdot C, N, D]$
};

\draw[->, thick, color=sorblue] (in_res) -- (fold);
\draw[->, thick, color=teal!80!black] (fold) -- (patching);
\draw[->, thick, color=teal!80!black] (patching) -- (intra);
\draw[->, thick, color=sororange!90!black] (intra) -- (pool);

% =========================================================================
% STAGE 2: DCT-II Spectral Conditioning Branch (x = 5.5)
% =========================================================================
\node[font=\bfseries\normalsize, text=violet!85!black, anchor=center] (s2_title) at (5.5, 0.4) {2. Spectral FiLM Branch};

\node[spec, below=0.22cm of s2_title] (dct) {
    \textbf{Orthonormal DCT-II}\\
    \footnotesize $\mathbf{C}_{\text{dct}} = \mathbf{X}_c \mathbf{D}^T \in \mathbb{R}^{B \cdot C \times L}$ (Energy compaction)
};

\node[spec, below=0.22cm of dct] (spec_power) {
    \textbf{Log-Power Subbands}\\
    \footnotesize $\mathbf{S} = \log(1 + \mathbf{C}_{\text{dct}}^2) \to \text{Pool}(K=4) \implies \bar{\mathbf{s}} \in [B \cdot C, K]$
};

\node[spec, below=0.22cm of spec_power] (film_mlp) {
    \textbf{Zero-Init FiLM Generator}\\
    \footnotesize MLP: $K \to \text{hidden} \to 2D \implies (\boldsymbol{\gamma}, \boldsymbol{\beta}) = (\mathbf{0}, \mathbf{0})$
};

\node[film_node, below=0.28cm of film_mlp] (film) {
    \textbf{FiLM Spectral Modulation}\\
    $\mathbf{H}_{\text{cond}} = \mathbf{H} \odot (1 + \boldsymbol{\gamma}) + \boldsymbol{\beta}$\\
    \footnotesize \textit{Step 0: Identity Initialization ($\mathbf{H}_{\text{cond}} = \mathbf{H}$)}
};

% Connectors in Stage 2
\draw[->, thick, color=violet!80!black] (dct) -- (spec_power);
\draw[->, thick, color=violet!80!black] (spec_power) -- (film_mlp);
\draw[->, thick, color=purple!90!black] (film_mlp) -- (film)
    node[midway, right, font=\footnotesize] {$(\boldsymbol{\gamma}, \boldsymbol{\beta})$};

% =========================================================================
% STAGE 3: Inter-Patch Mixer & Forecasting Head (x = 11.0)
% =========================================================================
\node[font=\bfseries\normalsize, text=sororange!95!black, anchor=center] (s3_title) at (11.0, 0.4) {3. Inter-Patch \& Forecast};

\node[inter, below=0.22cm of s3_title] (inter) {
    \textbf{Neighbor-Gated Inter-Mixer}\\
    \footnotesize Context: $[\mathbf{h}_{i-1}, \mathbf{h}_i, \mathbf{h}_{i+1}] \in 3D \to \text{MLP}_{3D \to D}$\\
    \footnotesize Dynamic Gate $g_i \in [0, 1]$ + Feat-MLP
};

\node[head, below=0.22cm of inter] (head) {
    \textbf{Forecasting Head}\\
    \footnotesize $\text{Linear}(N \cdot D \to D) \to \text{GELU} \to \text{Linear}(D \to H)$
};

\node[output, below=0.22cm of head] (out_res) {
    \textbf{Residual Forecast} $\hat{\mathbf{Y}}_{\text{res}} \in \mathbb{R}^{B \times H \times C}$\\
    \footnotesize Unfold channels + Periodic forecast $\hat{\mathbf{Y}}_{\text{cycle}}$
};

\node[card, below=0.22cm of out_res, text width=4.6cm] (summary_card) {
    \textbf{Key SORMixer Specifications:}\\
    $\bullet$ \textbf{Params}: $\sim \mathbf{100\text{k}}$ ($10\times$ smaller than PatchTST)\\
    $\bullet$ \textbf{Complexity}: $\mathcal{O}(N)$ strictly linear in patches\\
    $\bullet$ \textbf{Channel}: Independent residual modeling\\
    $\bullet$ \textbf{Conditioning}: Orthonormal DCT-II FiLM
};

% Connectors between Stages
% Branch from Stage 1 Fold to Stage 2 DCT
\draw[->, thick, color=violet!80!black] (fold.east) -- ++(0.45, 0) node[midway, above, font=\footnotesize] {$\mathbf{X}_c$} |- (dct.west);

% Stream from Stage 1 Pool to Stage 2 FiLM
\draw[->, thick, color=teal!80!black] (pool.east) -- ++(0.45, 0) node[midway, above, font=\footnotesize] {$\mathbf{H}$} |- (film.west);

% Stream from Stage 2 FiLM to Stage 3 Inter-Patch (centered in corridor)
\draw[->, thick, color=purple!80!black] (film.east) -- ++(0.45, 0) node[midway, above, font=\footnotesize] {$\mathbf{H}_{\text{cond}}$} |- (inter.west);

\draw[->, thick, color=sororange!90!black] (inter) -- (head);
\draw[->, thick, color=sorgreen!80!black] (head) -- (out_res);

\end{tikzpicture}
}%
}
\end{frame}

%######################################################################
% PART 3 --- LITERATURE: CycleNet, HDMixer, PatchTST  (TO BE ADDED)
%######################################################################
% >>> INSERT the "Literature" frames here. <<<

%######################################################################
% PART 4 --- WHICH PARAMETERS DO WHAT?
%######################################################################

\begin{frame}{General model parameters}
\scriptsize
\setlength{\tabcolsep}{5pt}
\renewcommand{\arraystretch}{1.1}
\begin{tabular}{@{}L{2.1cm} L{5.6cm} L{5.6cm}@{}}
  \toprule
  \textbf{Parameter} & \textbf{What it does and why it matters}
                     & \textbf{Relation to data} \\
  \midrule
  \texttt{SEQ\_LEN}
    & Length of the look-back window. Sets how many cycles the model
      sees at once. A window shorter than the dominant period prevents
      the recurrent-cycle component from learning a full cycle, and
      also caps the period the FFT can detect.
    & Should be a multiple of the dominant period. Long-cycle data
      needs larger windows; short-range noisy data can use smaller
      ones. \textbf{Most impactful parameter overall.} \\
  \addlinespace
  \texttt{PRED\_LENS}
    & Forecast horizons evaluated. Defines how far into the future
      the model must predict. Longer horizons are harder because the
      model relies more on learned periodic structure and less on
      recent observations.
    & A task specification, not a data property. Should cover
      multiple horizons to separate a model that generalises from one
      that happens to fit a single output length. \\
  \bottomrule
\end{tabular}
\end{frame}

\begin{frame}{SOR backbone parameters}
\scriptsize
\setlength{\tabcolsep}{5pt}
\renewcommand{\arraystretch}{1.1}
\begin{tabular}{@{}L{2.1cm} L{5.6cm} L{5.6cm}@{}}
  \toprule
  \textbf{Parameter} & \textbf{What it does and why it matters}
                     & \textbf{Relation to data} \\
  \midrule
  \texttt{SOR\_DROPOUT}
    & Dropout inside every mixer block and the head. The primary
      regularizer of the SOR backbone. Prevents memorising training
      windows and forces the residual forecast to generalise. Too
      high and the model cannot fit even the residual.
    & Tracks noise level and training-set size. Noisy data or small
      training splits need more dropout; clean data tolerates less.
      \textbf{Most impactful SOR-specific parameter.} \\
  \addlinespace
  \texttt{STRIDE\_RATIO}
    & Step between patches as a fraction of \texttt{PATCH\_LEN};
      controls overlap. Low values give denser coverage and stronger
      cross-patch context; high values give fewer, independent
      patches. Head size scales with the number of patches.
    & Denser overlap helps when adjacent patches are semantically
      related (periodic, smooth). \\
  \addlinespace
  \texttt{PATCH\_LEN}
    & Timesteps per patch; determines token granularity. Short
      patches preserve fine detail, long patches aggregate context.
      Ideal length aligns with the dominant period so each patch
      covers one coherent cycle phase.
    & Works best when it divides the dominant period cleanly. Short
      patches suit spiky signals; long patches suit smooth,
      single-cycle signals. \\
  \bottomrule
\end{tabular}
\end{frame}

%######################################################################
% PART 5 --- EXPERIMENTS
%######################################################################

\begin{frame}{Results at look-back $L=336$ (average over horizons)}
\centering
\fontsize{6.5}{8}\selectfont
\setlength{\tabcolsep}{3pt}
\renewcommand{\arraystretch}{1.2}
\begin{tabular}{@{}l | cc | cc | cc | cc || cc | cc@{}}
  \toprule
  \multirow{2}{*}{\textbf{Dataset}}
    & \multicolumn{2}{c}{\textbf{Linear}} & \multicolumn{2}{c}{\textbf{MLP}}
    & \multicolumn{2}{c}{\textbf{CNN}}    & \multicolumn{2}{c}{\textbf{SOR}}
    & \multicolumn{2}{c}{\textbf{PatchTST/42}} & \multicolumn{2}{c}{\textbf{HDMixer}} \\
    & MSE & MAE & MSE & MAE & MSE & MAE & MSE & MAE & MSE & MAE & MSE & MAE \\
  \midrule
  ETTh1   & 0.4164 & 0.4271 & 0.4527 & 0.4477 & 0.4436 & 0.4551 & 0.4327 & 0.4420 & 0.4155 & 0.4294 & \textbf{0.406} & \textbf{0.424} \\
  ETTh2   & 0.3601 & 0.4027 & 0.3738 & 0.4109 & 0.4025 & 0.4373 & 0.3646 & 0.4048 & \underline{0.3302} & \underline{0.3797} & \textbf{0.319} & \textbf{0.371} \\
  ETTm1   & 0.3594 & 0.3858 & 0.3597 & 0.3877 & 0.3967 & 0.4239 & 0.3659 & 0.3926 & \textbf{0.3518} & \underline{0.3817} & \underline{0.3525} & \textbf{0.3768} \\
  ETTm2   & 0.2668 & 0.3215 & 0.2768 & 0.3306 & 0.3098 & 0.3575 & 0.2790 & 0.3318 & \underline{0.2573} & \underline{0.3151} & \textbf{0.2513} & \textbf{0.3135} \\
  Weather & 0.2339 & 0.2722 & 0.2311 & 0.2699 & 0.2860 & 0.3171 & 0.2357 & 0.2788 & \underline{0.2297} & \underline{0.2657} & \textbf{0.2233} & \textbf{0.2635} \\
  Traffic & 0.4472 & 0.3209 & 0.4404 & 0.3106 & 0.5942 & 0.3357 & 0.4606 & 0.3100 & \textbf{0.3976} & \textbf{0.2669} & \underline{0.3988} & \underline{0.2695} \\
  Elec.   & 0.1754 & 0.2738 & 0.1695 & 0.2684 & 0.2377 & 0.3248 & 0.1829 & 0.2826 & \underline{0.1616} & \underline{0.2538} & \textbf{0.1593} & \textbf{0.2505} \\
  Solar   & \underline{0.2414} & \underline{0.2795} & \underline{0.2345} & \textbf{0.2754} & 0.2621 & 0.3000 & \textbf{0.2156} & 0.2860 & -- & -- & -- & -- \\
  \bottomrule
\end{tabular}

\vfill
{\fontsize{6}{7.2}\selectfont\textcolor{muted}{\textit{All models use $L=336$. \textbf{Bold} = best and
\underline{underline} = second best within each row, over the models
present. PatchTST/42 and HDMixer do not report Solar at $L=336$.
Our backbones: \texttt{seq\_len=336, batch=32, lr=1e-3}.}}}
\end{frame}

\begin{frame}{Hyperparameter tuning}
\small
\textit{Tuning shows whether the edge of SOR is architectural or just a lucky default.}\par
\vspace{0.5em}
{\setlength{\fboxsep}{7pt}\setlength{\fboxrule}{0.45pt}%
\noindent\color{black}\fbox{\begin{minipage}{\dimexpr\linewidth-2\fboxsep-2\fboxrule\relax}
\textbf{Why we tune}\par\vspace{0.25em}
The hand-picked configuration on the previous slide is a fair reference, but it is not optimal for any specific dataset.
\end{minipage}}}
\vspace{0.6em}

\textbf{What we tune, and why}\par\vspace{0.3em}
{\footnotesize
\renewcommand{\arraystretch}{1.12}
\begin{tabular}{@{}L{2.6cm} L{9.8cm}@{}}
  \texttt{seq\_len}      & History length; more cycles versus noise and cost. \\
  \texttt{patch\_len}    & Patch granularity; should align with the dominant period. \\
  \texttt{stride\_ratio} & Patch overlap; sets token count, receptive field and head size. \\
  \texttt{sor\_dropout}  & Main regularizer; balances under- and overfitting. \\
  \texttt{lr}            & Optimization stability and convergence within the tuning budget. \\
  \texttt{weight\_decay} & Global regularization; interacts with dropout and dataset size. \\
\end{tabular}}
\vspace{0.6em}

\textbf{Result on ETTh2: 201-trial random search}\par\vspace{0.3em}
{\small The best tuned SOR reduces test MSE by about \textbf{2.9\%} and MAE by about \textbf{0.9\%}. The gain survives fresh-seed validation.}
\end{frame}

%######################################################################
% PART 6 --- DISCUSSION, OUTLOOK AND CLOSING
%######################################################################

\begin{frame}{What did not work: the current CNN}
\small
\textit{Added convolutional processing did not improve the forecasts.}\par
\vspace{0.55em}
{\setlength{\fboxsep}{7pt}\setlength{\fboxrule}{0.45pt}%
\noindent\color{black}\fbox{\begin{minipage}{\dimexpr\linewidth-2\fboxsep-2\fboxrule\relax}
\textbf{Our interpretation}\par\vspace{0.25em}
The current CNN has the highest average MSE among our backbones on each dataset. We would deprioritize this configuration in the next round.
\end{minipage}}}
\vspace{0.75em}
\begin{itemize}
  \setlength{\itemsep}{0.55em}
  \item \textbf{Limited temporal context:} Two kernel-5 layers see nine samples locally. Interpolation resizes the output without learning a global history-to-future mapping.
  \item \textbf{Extra processing:} Hidden feature maps add computation and memory relative to Linear, even though this CNN has fewer parameters.
  \item \textbf{Possible overfitting:} Nonlinear filters may learn local noise after cycle removal, improving training fit without improving unseen forecasts.
\end{itemize}
\vspace{0.35em}
{\footnotesize\textcolor{muted}{Overfitting check: falling training loss with rising or persistently high validation loss. Test MSE alone does not establish the cause.}}
\end{frame}

\begin{frame}{What worked better, and why}
\small
\textit{The useful complexity depends on what remains after cycle removal.}\par
\vspace{0.55em}
{\setlength{\fboxsep}{7pt}\setlength{\fboxrule}{0.45pt}%
\noindent\color{black}\fbox{\begin{minipage}{\dimexpr\linewidth-2\fboxsep-2\fboxrule\relax}
\textbf{A stronger backbone needs to improve the residual forecast}\par\vspace{0.25em}
Cycle tables already handle recurring patterns. The remaining task may benefit more from the right temporal structure than from additional nonlinear layers.
\end{minipage}}}
\vspace{0.75em}
\begin{itemize}
  \setlength{\itemsep}{0.65em}
  \item \textbf{Linear on ETT:} Its direct mapping uses the full input window. If cycle removal leaves a simple residual, extra nonlinear processing may add little value.
  \item \textbf{MLP on Traffic and Electricity:} It also uses the full window, with nonlinear interactions that may capture more complex residual behavior.
  \item \textbf{SOR on Solar:} Patch mixing and DCT conditioning may help track changes in local shape and frequency content.
\end{itemize}
\vspace{0.35em}
{\footnotesize\textcolor{muted}{These are explanations to test with ablations. The scores show which configuration worked, but do not identify the reason on their own.}}
\end{frame}

\begin{frame}{What comes next?}
\large
{\small\textit{Explore combinations and a longer input window.}}\par
\vspace{1.35em}
\textcolor{accent}{\textbf{1}}\quad \textbf{Test backbone combinations}
\par\vspace{0.45em}
{\normalsize Combine MLP and CNN to test whether local features improve the global residual forecast.}

\vspace{1.5em}
\textcolor{accent}{\textbf{2}}\quad \textbf{Try sequence length 512}
\par\vspace{0.45em}
{\normalsize Compare $L=512$ with the current $L=336$ to assess the value of longer history.}

\vspace{1.25em}
{\footnotesize\textcolor{muted}{Keep splits and forecast horizons fixed. Compare validation MSE/MAE, runtime and memory before selecting the final model.}}
\end{frame}

\begin{frame}{Conclusion and questions}
\normalsize
\textit{The best forecasting design depends on the patterns left after cycle removal.}\par
\vspace{1.0em}
\begin{itemize}
  \setlength{\itemsep}{0.85em}
  \item \textbf{No single backbone leads everywhere.} Linear, MLP and SOR each performed best on different datasets in our comparison.
  \item \textbf{More complexity does not guarantee better forecasts.}
\end{itemize}
\vspace{1.1em}
{\setlength{\fboxsep}{9pt}\setlength{\fboxrule}{0.45pt}%
\noindent\color{black}\fbox{\begin{minipage}{\dimexpr\linewidth-2\fboxsep-2\fboxrule\relax}
\centering\large\textbf{Thank you. Any questions?}
\end{minipage}}}
\end{frame}

\end{document}
